\documentclass[twoside,11pt]{article}

\usepackage{jmlr2e}
\usepackage{amsmath}
\usepackage{microtype}
\usepackage{booktabs}
\usepackage{tabularx}
\usepackage{array}
\usepackage{enumitem}
\usepackage{lastpage}

\jmlrheading{0}{2026}{1-\pageref{LastPage}}{Submitted 2/2026}{Under Review}{PAPER-A-0001}{Carlos Jonnabio}
\ShortHeadings{XAI Benchmarking Framework and Operation Protocol}{XAI Benchmarking Framework}
\firstpageno{1}

\begin{document}

\title{A Framework for Rigorous Evaluation of Model-Agnostic Explainability Methods:\\
Multi-Metric Statistical Benchmarking, Operational Protocol, and Reproducibility}

\author{
\name Jonathan Herrera-Vasquez \email jonnabio@gmail.com \\
\addr UNADE
}

\editor{TBD}

\maketitle

\begin{abstract}
Evaluating explainability methods requires more than a single faithfulness proxy. We present a modular benchmarking framework centered on quantitative XAI quality metrics: fidelity, stability, sparsity, computational cost, and faithfulness gap, plus an explicit method for operating the framework end-to-end. On the UCI Adult benchmark, we use a staged protocol (EXP1 calibration/reproducibility and EXP2 comparative/robustness benchmarking); the robustness cohort currently contains 250 of 300 planned configurations (83.3\% coverage). Across complete model-size blocks ($5$ models, $N \in \{50,100,200\}$), Friedman tests indicate significant method differences for fidelity ($\chi^2=42.12$, $p=3.78\times10^{-9}$), stability ($\chi^2=43.88$, $p=1.60\times10^{-9}$), sparsity ($\chi^2=35.64$, $p=8.92\times10^{-8}$), faithfulness gap ($\chi^2=45.00$, $p=9.25\times10^{-10}$), and runtime ($\chi^2=27.72$, $p=4.16\times10^{-6}$). SHAP leads on fidelity/stability, DiCE leads on sparsity, and LIME is generally fastest and most practical outside SVM-KernelSHAP bottlenecks. We release the framework, operation protocol, and artifacts with explicit data-quality caveats for reproducible benchmark use. LLM-based semantic scoring is deferred to a dedicated follow-up study.
\end{abstract}

\begin{keywords}
Explainable AI, benchmark methodology, reproducibility, statistical evaluation, model-agnostic explanations
\end{keywords}

\section{Introduction}
XAI evaluation remains fragmented across incompatible metrics and ad hoc setups. Most comparisons over-index on fidelity and omit stability, sparsity, and semantic interpretability. This work operationalizes a benchmark-first approach:
\begin{enumerate}[leftmargin=*]
    \item A generic trainer/explainer architecture for model-agnostic evaluation.
    \item A prescriptive operation method that defines how the framework is run, audited, and reported.
    \item Multi-metric quantitative scoring with explicit trade-off analysis.
    \item Reproducibility and statistical rigor suitable for benchmark publication.
\end{enumerate}

\section{Related Work and Framework Design}

\subsection{Literature-Grounded Rationale}
The literature motivating this design supports five requirements:
\begin{enumerate}[leftmargin=*]
    \item Multi-metric evaluation is necessary.
    \item Faithfulness is multi-faceted.
    \item Computational tractability must be reported explicitly.
    \item Benchmark transparency and artifact accounting are first-order requirements.
    \item LLM-judge scoring should be treated as a separate calibrated track.
\end{enumerate}

\subsection{Architecture}
The framework decouples model training, explanation generation, and evaluation through a registry/factory design. Current model coverage includes \texttt{logreg}, \texttt{rf}, \texttt{xgb}, \texttt{svm}, and \texttt{mlp}, with \texttt{lime}, \texttt{shap}, \texttt{anchors}, and \texttt{dice} explainers.

\subsection{Metric Suite (Primary)}
Primary quantitative metrics are:
\begin{itemize}[leftmargin=*]
    \item Fidelity
    \item Stability
    \item Sparsity
    \item Cost
    \item Faithfulness gap
\end{itemize}

\subsection{Paper A Operational Outcome: Framework Operation Method (FOM-7)}
Paper A contributes not only benchmark findings, but also a reusable operating method for end-to-end benchmark execution, auditing, and claim formation.

\section{Methodology}

\subsection{Experimental Cohorts and Scope}
We intentionally organize evidence into staged cohorts so that development, reliability qualification, and confirmatory inference are not conflated. This design reduces circularity (tuning and claiming on the same artifacts) and makes the evidence path for each claim auditable.

We separate evidence into four cohorts:
\begin{itemize}[leftmargin=*]
    \item \textbf{EXP1 pilot/calibration} (\texttt{exp1\_adult/results})
    \item \textbf{EXP1 reproducibility cohort} (\texttt{exp1\_adult/reproducibility})
    \item \textbf{EXP2 comparative} (\texttt{exp2\_comparative/results})
    \item \textbf{EXP2 robustness cohort} (\texttt{exp2\_scaled/results})
\end{itemize}

\begin{table}[t]
\centering
\caption{Role of each evidence cohort in Paper A.}
\begin{tabularx}{\linewidth}{l X X}
\toprule
\textbf{Cohort} & \textbf{Primary objective} & \textbf{Role in manuscript claims} \\
\midrule
EXP1 pilot/calibration & Verify metric behavior and tune explainer/evaluation settings before large-scale execution. & Used for implementation sanity and calibration only; not used for primary significance claims. \\
EXP1 reproducibility & Measure run-to-run variance under repeated seeds on core configurations. & Used for reproducibility evidence (CV and dispersion profiles). \\
EXP2 comparative & Provide complete method/model breadth under fixed-seed constraints. & Used as balanced comparative layer and coverage check. \\
EXP2 robustness & Provide multi-seed and multi-sample-size evidence for inferential testing under controlled heterogeneity. & Primary inferential cohort for Friedman/Nemenyi/Wilcoxon analyses and cross-method conclusions. \\
\bottomrule
\end{tabularx}
\label{tab:cohorts}
\end{table}

This EXP1/EXP2 separation is methodologically necessary for four reasons: (i) it preserves tuning-evaluation separation, (ii) it isolates implementation validity from comparative validity, (iii) it keeps artifact failures explicit instead of hidden in pooled aggregates, and (iv) it constrains claims to cohorts that are fit for each claim type.

Operationally, the structure is analogous to train/dev/test discipline, adapted to XAI system evaluation: EXP1 functions as development and reliability qualification, while EXP2 functions as comparative and inferential qualification.

Accordingly, Paper A inferential claims are driven by EXP2, while EXP1 is retained for calibration and reproducibility reporting.

\subsection{Research Questions and Hypotheses}
This section specifies the inferential targets and separates confirmatory
discrimination questions from robustness/stability questions.

Let $\mathcal{K}$ denote explanation methods, let $\mathcal{B}$ denote
evaluation blocks with $b=(\text{model},N)$, and let $y_{k,b,m}$ denote the
run-level aggregated score for method $k$ in block $b$ on metric $m$.
We evaluate
$m \in \{\text{fidelity}, \text{stability}, \text{sparsity},
\text{faithfulness gap}, \text{cost}\}$.
These metrics represent distinct quality dimensions; therefore,
single-metric dominance is not interpreted as global superiority.

\paragraph{RQ1 (method separation).}
Do explanation methods differ significantly on each metric when compared over
repeated blocks?
For each metric $m$:
\begin{align}
H_{0,m}^{(1)} &: \text{all methods have equal median rank across blocks},\\
H_{1,m}^{(1)} &: \text{at least one method differs in median rank.}
\end{align}
Rejection of $H_{0,m}^{(1)}$ indicates that the benchmark discriminates methods
on metric $m$ under observed variance and coverage.

\paragraph{RQ2 (trade-off stability).}
Are method trade-offs stable across model families and sample sizes, or
context-dependent?
We formalize this through two complementary hypotheses:
\begin{align}
H_{0,m}^{(2a)} &: \text{method ordering for metric } m \text{ is stable across blocks},\\
H_{1,m}^{(2a)} &: \text{method ordering for metric } m \text{ shifts across blocks},\\
H_{0,m}^{(2b)} &: \text{pairwise method differences on } m \text{ do not vary systematically}\\
              &\quad \text{with model family or sample size},\\
H_{1,m}^{(2b)} &: \text{pairwise method differences on } m \text{ vary systematically}\\
              &\quad \text{with model family and/or sample size.}
\end{align}

RQ2 is interpreted as a portability question: we treat context-dependent rank
reversals or large cross-stratum effect shifts as evidence that trade-off
guidance is not stable.

Methodological scoping in this manuscript is explicit: RQ1 is the primary
confirmatory layer; RQ2 is the robustness layer connecting inferential outcomes
to external validity and deployment relevance.

\subsection{Dataset, Preprocessing, and Leakage Controls}
The Adult Income benchmark is loaded through a dedicated pipeline
(\path{src/data_loading/adult.py}) with explicit source fallback,
schema control, and transformation traceability.

\paragraph{Data provenance and schema control.}
The loader uses OpenML Adult v2 (\texttt{data\_id=1590}) as primary source.
If unavailable, it falls back to direct UCI files
(\texttt{adult.data}, \texttt{adult.test}).
Cached artifacts are checksum-validated and stored with source and schema
metadata (shape, columns, timestamp, and package versions).
The enforced schema has 15 columns: one binary target (\texttt{income}),
6 numeric predictors, and 8 categorical predictors.

\begin{table}[t]
\centering
\caption{Adult dataset feature groups used in preprocessing.}
\begin{tabularx}{\linewidth}{l X}
\toprule
\textbf{Group} & \textbf{Variables} \\
\midrule
Numeric (6) & \texttt{age, fnlwgt, education-num, capital-gain, capital-loss, hours-per-week} \\
Categorical (8) & \texttt{workclass, education, marital-status, occupation, relationship, race, sex, native-country} \\
Target & \texttt{income} mapped to \{0,1\} \\
\bottomrule
\end{tabularx}
\label{tab:adult-schema}
\end{table}

Cleaning includes whitespace normalization, conversion of ``?'' markers to
\texttt{NaN}, target canonicalization
(\texttt{<=50K}/\texttt{<=50K.}$\rightarrow 0$,
\texttt{>50K}/\texttt{>50K.}$\rightarrow 1$), duplicate removal, and numeric
type coercion. In the current artifact state, the cleaned corpus contains
48{,}790 instances.

\paragraph{Split discipline and deterministic execution.}
Data are split using a stratified holdout protocol:
\texttt{train\_test\_split(test\_size=0.2, stratify=y, random\_state=seed)}.
This preserves class prevalence across partitions while enabling exact
reproducibility for each seed and controlled seed variation across cohorts.

\paragraph{Preprocessing design.}
Transformations are implemented with a \texttt{ColumnTransformer}:
\begin{itemize}[leftmargin=*]
    \item Numeric branch: \texttt{StandardScaler} on six numeric variables.
    \item Categorical branch: \texttt{OneHotEncoder} on eight categorical variables.
    \item Encoder settings: \texttt{handle\_unknown='ignore'} and \texttt{sparse\_output=False}.
\end{itemize}
The encoder is fitted on training data only. Unseen categories in test
instances are safely ignored, preventing transform-time failures. In the
current training vocabulary, one-hot expansion yields 102 categorical columns;
with 6 scaled numeric columns, the processed design matrix has 108 features.

\paragraph{Leakage controls and feature-space alignment.}
Leakage prevention is enforced by fitting transformations on training data only
and reusing the fitted transformer for test and explanation pipelines.
Feature names are extracted from the fitted transformer and propagated to
explainers and metric engines to avoid attribution-index mismatch. In
\texttt{ExperimentRunner}, when a persisted preprocessor artifact exists
(\texttt{preprocessor.pkl} or \texttt{preprocessor.joblib}), it is loaded and
reused so that explanation-time transformations exactly match model-training
transformations.

\paragraph{Construct caveat.}
Adult is suitable for tabular XAI benchmarking, but contains socio-demographic
attributes and proxies (e.g., \texttt{sex}, \texttt{race},
\texttt{native-country}). The protocol here addresses preprocessing and leakage
rigor; it does not remove fairness-related construct limitations of the data.

\subsection{Experimental Design and Analysis Units}
\subsubsection{Design Type and Inferential Target}
The EXP2 comparative study is a crossed, repeated-seed factorial design with
model family $g \in \mathcal{G}$, explainer $k \in \mathcal{K}$, seed
$s \in \mathcal{S}$, and sampling intensity $n \in \mathcal{N}$, where
\[
\mathcal{G}=\{\texttt{logreg},\texttt{rf},\texttt{xgb},\texttt{svm},\texttt{mlp}\},
\]
\[
\mathcal{K}=\{\texttt{lime},\texttt{shap},\texttt{anchors},\texttt{dice}\},
\]
\[
\mathcal{S}=\{42,123,456,789,999\}, \qquad
\mathcal{N}=\{50,100,200\}.
\]
Each executed configuration yields a run-level value $y_{g,k,s,n,m}$ for metric
$m$. Because seed coverage is partially unbalanced in realized artifacts,
confirmatory inference is defined on block-level summaries over
$b=(g,n)$:
\[
\bar{y}_{k,b,m}=\frac{1}{|S_{k,b}|}\sum_{s \in S_{k,b}} y_{g,k,s,n,m},
\]
where $S_{k,b}\subseteq\mathcal{S}$ denotes analyzable seeds for method $k$ in
block $b$. This prevents methods with more successful seed executions from
receiving disproportionate influence in omnibus tests.

\subsubsection{Planned Matrix and Realized Coverage}
The planned robustness matrix contains
\[
|\mathcal{G}|\times|\mathcal{K}|\times|\mathcal{S}|\times|\mathcal{N}|
=5\times4\times5\times3=300
\]
run configurations (see
\path{configs/experiments/exp2_scaled/manifest.yaml}). In the current artifact
snapshot, 250 files are present in
\path{experiments/exp2_scaled/results}, of which 233 are analyzable,
16 are empty, and 1 is malformed JSON.

For global method comparisons, we retain only block-complete model-size
contexts with at least one analyzable run per method. This yields $15/15$
complete blocks ($5$ models $\times$ $3$ sample sizes). Within those blocks,
seed availability is heterogeneous:
\begin{itemize}[leftmargin=*]
    \item SHAP: 2--5 seeds per block (mean $\approx 4.33$)
    \item LIME: 5 seeds per block (mean $=5.00$)
    \item Anchors: 1--5 seeds per block (mean $\approx 2.87$)
    \item DiCE: 2--5 seeds per block (mean $\approx 3.33$)
\end{itemize}
No imputation is applied for missing runs.

\subsubsection{Within-Run Sampling Protocol}
Each run uses stratified error-quadrant sampling
(\path{src/evaluation/sampler.py}) over $\{\mathrm{TP},\mathrm{TN},\mathrm{FP},\mathrm{FN}\}$.
The configuration parameter \texttt{samples\_per\_class} is interpreted as
target samples per quadrant, so nominal run size is $4N$.

Realized run size is lower when (i) one or more quadrants has fewer than $N$
available test instances for a given model/seed or (ii) sampled instances fail
during explainer evaluation and are dropped. Across analyzable EXP2 runs,
realized run size ranges from 27 to 800 instances (median 400). Across all
retained instance evaluations (103,577 total), quadrant composition remains
near-balanced: TP 26,523 (25.6\%), TN 26,050 (25.2\%), FP 25,200 (24.3\%), and
FN 25,804 (24.9\%).

\subsubsection{Analysis Units and Anti-Pseudoreplication Strategy}
The design separates instance-level logging from inferential units:
\begin{enumerate}[leftmargin=*]
    \item Instance level: raw per-case metric outputs.
    \item Run level: per-configuration aggregate $y_{g,k,s,n,m}$.
    \item Block level: per-method summaries $\bar{y}_{k,b,m}$ for $b=(g,n)$.
    \item Matched-pair level: SHAP--LIME contrasts on shared $(g,s,n)$ cells.
\end{enumerate}

Global Friedman/Nemenyi tests are run on a $15\times4$ block-method matrix
(15 model-size blocks, 4 methods), avoiding pseudo-replication from
instance-level values and preventing seed replicates from being treated as
independent blocks. For SHAP--LIME contrasts, matched analyzable sets are:
\[
\left|\mathcal{P}_{\text{all models}}\right|=65,\qquad
\left|\mathcal{P}_{\texttt{logreg/rf/xgb}}\right|=45.
\]

\subsubsection{Retention and Exclusion Rules}
Inferential datasets are constructed under fixed artifact-qualification rules:
\begin{enumerate}[leftmargin=*]
    \item Parseable JSON artifact required (malformed files excluded).
    \item Non-empty \texttt{instance\_evaluations} required for run-level analysis.
    \item No synthetic completion of missing runs or seeds.
    \item Omnibus tests restricted to block-complete contexts.
    \item Paired tests restricted to explicitly matched SHAP--LIME cells.
\end{enumerate}
This makes the distinction explicit between design intent (300-run full
factorial) and inferential realization (artifact-qualified subsets).

\subsection{Experimental Design Protocol (Replication Specification)}
This subsection specifies the executable protocol used for Paper A. The goal is
to make the EXP2 pipeline reproducible without undocumented assumptions.

\subsubsection{Purpose and Hypotheses Linkage}
The protocol instantiates Section~3.2 questions in operational form:
\begin{enumerate}[leftmargin=*]
    \item RQ1: method separation across primary quantitative metrics.
    \item RQ2: stability of method trade-offs across model family and sampling intensity.
\end{enumerate}
Reliability and faithfulness are evaluated as properties of the full
$(\text{model},\text{explainer},\text{data},\text{metric})$ system.

\subsubsection{Experimental Factors and Conditions}
Independent variables are defined by
\path{configs/experiments/exp2_comparative} and
\path{configs/experiments/exp2_scaled}:
\begin{enumerate}[leftmargin=*]
    \item Dataset (fixed): Adult Income (\texttt{dataset: adult}).
    \item Model family: \texttt{logreg}, \texttt{rf}, \texttt{xgb}, \texttt{svm}, \texttt{mlp}.
    \item Explainer method: \texttt{lime}, \texttt{shap}, \texttt{anchors}, \texttt{dice}.
    \item Seeds (robustness cohort): $\{42,123,456,789,999\}$.
    \item Sampling intensity: \texttt{samples\_per\_class} $= N$, $N\in\{50,100,200\}$.
    \item Stability perturbation regime: Gaussian noise with
    \texttt{stability\_perturbations}=15 and default
    \texttt{stability\_noise\_level}=0.1.
    \item Method-specific secondary flag: \texttt{counterfactual=true} for DiCE.
\end{enumerate}
Planned robustness design size is $5\times4\times5\times3=300$ configurations
(see Table X).

Ablation scope is explicitly fixed for Paper A:
\begin{enumerate}[leftmargin=*]
    \item Paper A is a comparative benchmark, not a component-ablation study.
    \item No component-removal or hyperparameter-factor ablation enters confirmatory inference.
    \item EXP1 tuning artifacts are calibration records only and are excluded from inferential sets.
    \item Ablation-driven causal attribution is deferred to a separate extension study (Paper B).
\end{enumerate}

\subsubsection{Controls and Baselines}
Within each matched cell $(\text{model},\text{seed},N)$, the following are held constant:
\begin{enumerate}[leftmargin=*]
    \item Black-box model artifact path.
    \item Preprocessor artifact and transformed feature space.
    \item Split and sampling seeds.
    \item Evaluation sampler and metric engine code path.
    \item Primary metric set and run-level aggregation rules.
\end{enumerate}
Comparator set is \texttt{lime/shap/anchors/dice}. Sanity-control baselines
(random attribution, label randomization, parameter randomization) are not yet
included.
Therefore, Paper A claims are restricted to relative method performance under
trained-model conditions; no claim is made about explanation validity under
label/model randomization perturbations.

\subsubsection{Data Protocol}
Data handling follows Section~3.3 and \path{src/data_loading/adult.py}:
\begin{enumerate}[leftmargin=*]
    \item Adult cleaning: whitespace normalization, missing-value canonicalization,
    target canonicalization, duplicate removal.
    \item Stratified split:
    \texttt{train\_test\_split(test\_size=0.2, stratify=y, random\_state=seed)}.
    \item Preprocessor fit on training split only (unless persisted preprocessor is provided).
    \item Shared train/test transformation and feature-name propagation.
    \item TP/TN/FP/FN stratified evaluation sampling with target $N$ per quadrant.
\end{enumerate}
For seed 42, split sizes are train 39,032 and test 9,758
($y=0$: 29,687/7,422; $y=1$: 9,345/2,336). Nominal run size is $4N$, with
lower realized size when quadrant availability is below $N$.

\subsubsection{Model Training Protocol}
EXP2 loads frozen EXP1 model artifacts; models are not retrained during EXP2.
Binary target objective is income classification.

Persisted model configurations:
\begin{enumerate}[leftmargin=*]
    \item LogReg: \texttt{solver=lbfgs}, \texttt{max\_iter=1000}, \texttt{random\_state=42}.
    \item RF: \texttt{n\_estimators=50}, \texttt{max\_depth=15},
    \texttt{min\_samples\_leaf=1}, \texttt{class\_weight=balanced\_subsample},
    \texttt{random\_state=42}.
    \item XGB: \texttt{n\_estimators=100}, \texttt{max\_depth=6},
    \texttt{learning\_rate=0.1}, \texttt{objective=binary:logistic},
    \texttt{eval\_metric=logloss}, \texttt{random\_state=42}.
    \item SVM: \texttt{kernel=rbf}, \texttt{C=1.0}, \texttt{gamma=scale},
    \texttt{probability=true}, \texttt{random\_state=42}.
    \item MLP: \texttt{hidden\_layer\_sizes=[100]}, \texttt{activation=relu},
    \texttt{solver=adam}, \texttt{max\_iter=500}, \texttt{random\_state=42}.
\end{enumerate}
Early stopping is implemented only in \texttt{XGBoostTrainer} when validation
data are passed; it is not active in the released EXP2 model artifacts.
No post-hoc probability calibration is applied.
Calibration interpretation policy:
\begin{enumerate}[leftmargin=*]
    \item All methods are evaluated on identical uncalibrated model outputs per matched cell.
    \item Inferential statements are comparative (ranking/trade-off), not calibration claims.
    \item No calibration-dependent claim is made in Paper A.
\end{enumerate}

\subsubsection{Explanation Generation Protocol}
All explainers are called through \texttt{ExplainerWrapper.explain\_instance}
and timed per instance.

\paragraph{SHAP.}
\texttt{tree} mode for RF/XGB and \texttt{kernel} mode for SVM/MLP/LogReg,
\texttt{n\_background\_samples=50}. Tree mode uses interventional perturbation and
probability output; class-1 attributions are retained.

\paragraph{LIME.}
\texttt{num\_samples=1000}, \texttt{num\_features=10},
\texttt{kernel\_width=3.0}, \texttt{discretize\_continuous=False};
class-1 explanation map is converted to dense vector.

\paragraph{Anchors.}
\texttt{AnchorTabular} fitted on transformed training data; rule threshold 0.95;
anchor membership mapped to binary feature-importance vector.

\paragraph{DiCE.}
Backend method \texttt{random}; \texttt{desired\_class="opposite"};
importance proxy is absolute original--counterfactual feature difference.

Implementation caveats:
\begin{enumerate}[leftmargin=*]
    \item Anchors threshold is hardcoded at 0.95 in wrapper call.
    \item DiCE YAML \texttt{num\_counterfactuals} is not propagated; runtime uses
    \texttt{total\_CFs=1}.
\end{enumerate}
Parameter-binding policy for reported runs:
\begin{enumerate}[leftmargin=*]
    \item SHAP: \texttt{explainer\_type} and \texttt{n\_background\_samples} from YAML are applied.
    \item LIME: \texttt{kernel\_width} from YAML params is applied; \texttt{num\_samples} and \texttt{num\_features} are bound through schema fields (defaults 1000/10), with duplicate \texttt{params} entries dropped.
    \item Anchors: YAML \texttt{threshold} is not consumed; effective threshold is fixed at 0.95.
    \item DiCE: YAML \texttt{num\_counterfactuals} is not consumed; effective setting is \texttt{total\_CFs=1} with \texttt{method=random}.
\end{enumerate}
These bindings define the effective treatment conditions for EXP2 claims.

Runtime constraints: explanation cost is measured with \texttt{perf\_counter};
no explicit timeout/memory cap in YAML; batch runner enforces one evaluation
worker per experiment process to prevent nested parallelism explosion.
Run-level cost profile over analyzable EXP2 runs
(\path{outputs/analysis/paper_a_exp2_stats/exp2_run_level_metrics.csv}):
\begin{center}
\begin{tabular}{lrrrrr}
\toprule
Method & Runs & Mean (ms) & Median (ms) & P95 (ms) & Max (ms) \\
\midrule
SHAP    & 65 & 503870.23 & 1098.11 & 1603727.03 & 9393990.21 \\
LIME    & 75 & 3660.68   & 65.73   & 13192.05   & 85825.65 \\
Anchors & 43 & 28762.43  & 34063.36 & 64568.75  & 67997.67 \\
DiCE    & 50 & 15337.17  & 9759.79 & 34172.81   & 37164.19 \\
\bottomrule
\end{tabular}
\end{center}

\subsubsection{Evaluation Metrics and Measurement Procedure}
Metrics are computed per instance in \texttt{MetricsEngine}:
\begin{enumerate}[leftmargin=*]
    \item Cost: explanation time in milliseconds.
    \item Sparsity: active-feature ratio ($|w_i|>\tau$, $\tau=10^{-4}$).
    \item Fidelity proxy: faithfulness correlation score.
    \item Faithfulness gap: prediction shift after masking top-$k$ features ($k=5$).
    \item Stability: mean pairwise cosine similarity under Gaussian perturbations.
\end{enumerate}

Faithfulness implementation (\path{src/metrics/faithfulness.py}):
\begin{align}
\Delta_k &= \left|f(x)-f\left(x_{\text{mask-top-}k}\right)\right|,\\
\rho &= \mathrm{corr}\left(|w_i|,\left|f(x)-f(x_{-i})\right|\right).
\end{align}
Pipeline mapping is \texttt{fidelity} $\leftarrow \rho$ and
\texttt{faithfulness\_gap} $\leftarrow \Delta_k$.

Stability implementation (\path{src/metrics/stability.py}):
\begin{align}
x^{(t)} &= x + \epsilon^{(t)},\qquad \epsilon^{(t)}\sim\mathcal{N}(0,\sigma^2),\\
S &= \frac{2}{T(T-1)}\sum_{a<b}\cos\left(e^{(a)},e^{(b)}\right).
\end{align}

Secondary track: \texttt{cf\_sensitivity} is computed when
\texttt{counterfactual=true}; domain-alignment metric exists but is disabled in
EXP2 configs (\texttt{domain=false}).

Aggregation is instance $\rightarrow$ run (\texttt{mean/std/min/max/count})
$\rightarrow$ block summary over $(\text{model},N)$.

EEU note: Explanation Energy Units are specified in metric-design artifacts but
not computed in current EXP2 runtime.
EEU is therefore excluded from confirmatory inference, and Paper A cost claims
are defined strictly on wall-clock milliseconds.

\subsubsection{Statistical Analysis Plan}
Confirmatory dataset construction follows artifact qualification rules (Sections
3.4.5 and 3.6). Planned robustness grid is 300 runs; current snapshot has 250
present and 233 analyzable runs.

Inferential procedures:
\begin{enumerate}[leftmargin=*]
    \item Friedman test per metric over methods with $(\text{model},N)$ blocks.
    \item Nemenyi post-hoc pairwise test when Friedman is significant.
    \item Paired Wilcoxon SHAP-vs-LIME on matched $(\text{model},\text{seed},N)$ cells.
\end{enumerate}
Nemenyi provides family-wise control for pairwise localization. Across-metric
multiplicity is controlled with Holm-Bonferroni adjustment over the five
primary metrics, applied separately to the Friedman family and to each Wilcoxon
family (45-pair and 65-pair matched sets). Uncertainty is reported with
mean/std/CV and, where generated, $t$-based and bootstrap CIs.

Effect-size reporting is exported with inference artifacts:
\begin{enumerate}[leftmargin=*]
    \item Friedman layer: Kendall's $W$ for each metric.
    \item Paired SHAP--LIME layer: Cohen's $d_z$ and median paired differences.
\end{enumerate}

\subsubsection{Robustness and Sensitivity Analyses}
Implemented robustness axes:
\begin{enumerate}[leftmargin=*]
    \item Seed heterogeneity (5 seeds).
    \item Sample-size heterogeneity ($N\in\{50,100,200\}$ per quadrant).
    \item Perturbation stability under Gaussian noise ($T=15$, $\sigma=0.1$ default).
\end{enumerate}
Additional reproducibility sensitivity is reported from the EXP1 repeated-seed
cohort (\path{reproducibility_report.csv}).

Not included in Paper A core protocol: adversarial perturbation campaigns and
causal-intervention benchmarks beyond proxy counterfactual sensitivity.
Robustness interpretation in Paper A is limited to seed, sample-size, and
Gaussian-noise perturbation dimensions.
Reporting rule for robustness consistency:
\begin{enumerate}[leftmargin=*]
    \item SHAP--LIME median-difference sign must match between 45-cell and 65-cell matched analyses.
    \item Holm-adjusted significance decision must not flip between these analyses.
\end{enumerate}
Metrics violating either condition are reported as sensitivity-unstable.

\subsubsection{Reproducibility Package and Rerun Instructions}
Release package currently includes source code, experiment YAMLs, model and
preprocessor artifacts, per-run JSON/CSV outputs, batch manifest, and figure
generation scripts.

Minimal rerun sequence:
\begin{enumerate}[leftmargin=*]
    \item Install dependencies from \texttt{requirements-frozen.txt}.
    \item Verify or regenerate \path{experiments/exp1_adult/models} artifacts.
    \item Run EXP2 robustness batch:
    \texttt{python3 scripts/run\_batch\_experiments.py --config-dir configs/experiments/exp2\_scaled --output outputs/batch\_results.csv --parallel}.
    \item Run coverage audit: \texttt{python3 scripts/check\_missing\_results.py}.
    \item Generate summaries and figures.
\end{enumerate}

Execution provenance anchor for this manuscript snapshot:
\begin{enumerate}[leftmargin=*]
    \item Reference regeneration host: Apple M3 Pro, macOS 26.3 (arm64), 11 logical CPUs, 18 GB RAM, Python 3.13.2.
    \item Dependency lock hash: \texttt{sha256(requirements-frozen.txt)} =\\
    \texttt{b52d5a1f2e2edd5ada372ca66d18ef1447712fdd2c21f004d7e1f46a5ef9c6dc}.
    \item Immutable DOI is not yet minted in this repository state; current provenance key is \path{outputs/batch_manifest.json} with git hash \texttt{9fc70eb1e218a11f2fdd4bdb3aab3ea10a799a6f}.
\end{enumerate}

\subsection{Statistical Analysis Pipeline (Replication-Grade Specification)}
This subsection defines the executable protocol used to produce all Paper A
inferential artifacts from repository-tracked inputs (see Table~X and Appendix~X).

\subsubsection{Experimental Factors and Conditions}
Independent variables are instantiated from
\path{configs/experiments/exp2_comparative},
\path{configs/experiments/exp2_scaled}, and
\path{configs/experiments/exp2_scaled/manifest.yaml}:
\begin{enumerate}[leftmargin=*]
    \item Dataset (fixed): \texttt{adult}.
    \item Model family: \texttt{logreg}, \texttt{rf}, \texttt{xgb}, \texttt{svm}, \texttt{mlp}.
    \item Explainer method: \texttt{shap}, \texttt{lime}, \texttt{anchors}, \texttt{dice}.
    \item Seeds: $\{42,123,456,789,999\}$.
    \item Sampling intensity: \texttt{samples\_per\_class}=N, $N\in\{50,100,200\}$ (robustness cohort).
    \item Stability perturbation regime: \texttt{stability\_perturbations}=15 and default \texttt{stability\_noise\_level}=0.1.
    \item Counterfactual metric activation: \texttt{counterfactual=true} only for DiCE runs.
\end{enumerate}
Planned robustness grid size is $5\times4\times5\times3=300$. No
component-removal ablation is included in Paper A confirmatory inference.

\subsubsection{Controls and Baselines}
Within each matched $(\text{model},\text{seed},N)$ cell, the following are held
constant:
\begin{enumerate}[leftmargin=*]
    \item Model artifact path.
    \item Preprocessor artifact and transformed feature space.
    \item Split/sampling random-state values.
    \item Evaluation sampler and metric-engine implementation.
    \item Metric set and aggregation logic.
\end{enumerate}
Baselines are cross-method (\texttt{shap/lime/anchors/dice}) under shared
conditions. Random-attribution, label-randomization, and
model-randomization sanity controls are not part of Paper A runs; therefore
claims are bounded to trained-model comparative performance.

\subsubsection{Data Protocol}
Data handling is implemented in \path{src/data_loading/adult.py} and
\path{src/evaluation/sampler.py}:
\begin{enumerate}[leftmargin=*]
    \item Adult cleaning (whitespace normalization, missing-value canonicalization, duplicate removal, target normalization).
    \item Stratified split with \texttt{test\_size=0.2}, \texttt{stratify=y}, \texttt{random\_state=seed}.
    \item Preprocessor fit on train only.
    \item Shared train/test transformation and feature-name propagation.
    \item TP/TN/FP/FN stratified sampling with target size $N$ per stratum.
\end{enumerate}
For seed 42, split sizes are train 39,032 and test 9,758
($y=0$: 29,687/7,422; $y=1$: 9,345/2,336). Nominal run size is $4N$; realized
size may be lower under stratum-capacity constraints.

\subsubsection{Model Training Protocol}
EXP2 evaluates frozen EXP1 model artifacts; no retraining occurs in EXP2.
Persisted model settings are:
\begin{enumerate}[leftmargin=*]
    \item LogReg: \texttt{solver=lbfgs}, \texttt{max\_iter=1000}, \texttt{random\_state=42}.
    \item RF: \texttt{n\_estimators=50}, \texttt{max\_depth=15}, \texttt{min\_samples\_leaf=1}, \texttt{class\_weight=balanced\_subsample}, \texttt{random\_state=42}.
    \item XGB: \texttt{n\_estimators=100}, \texttt{max\_depth=6}, \texttt{learning\_rate=0.1}, \texttt{objective=binary:logistic}, \texttt{eval\_metric=logloss}, \texttt{random\_state=42}.
    \item SVM: \texttt{kernel=rbf}, \texttt{C=1.0}, \texttt{gamma=scale}, \texttt{probability=true}, \texttt{random\_state=42}.
    \item MLP: \texttt{hidden\_layer\_sizes=[100]}, \texttt{activation=relu}, \texttt{solver=adam}, \texttt{max\_iter=500}, \texttt{random\_state=42}.
\end{enumerate}
Early stopping is implemented only for XGBoost when validation data are passed;
it is inactive in persisted EXP2 artifacts. No post-hoc probability calibration
is applied.

\subsubsection{Explanation Generation Protocol}
Explanations are generated via \texttt{ExplainerWrapper.explain\_instance} on the
dense transformed feature vector (108 dimensions in current Adult vocabulary).
Per-instance runtime is recorded via wall-clock timing.
\begin{enumerate}[leftmargin=*]
    \item SHAP: \texttt{tree} mode for RF/XGB, \texttt{kernel} for LogReg/SVM/MLP; \texttt{n\_background\_samples=50}; class-1 attributions retained.
    \item LIME: \texttt{num\_samples=1000}, \texttt{num\_features=10}, \texttt{kernel\_width=3.0}, \texttt{discretize\_continuous=False}.
    \item Anchors: \texttt{AnchorTabular} on transformed train matrix; effective precision threshold fixed at 0.95.
    \item DiCE: backend \texttt{method=random}, \texttt{desired\_class=opposite}, absolute original--counterfactual deltas as importance proxy.
\end{enumerate}
Effective parameter binding caveats:
\begin{enumerate}[leftmargin=*]
    \item Anchors YAML \texttt{threshold} is not consumed; runtime threshold is fixed at 0.95.
    \item DiCE YAML \texttt{num\_counterfactuals} is not consumed; runtime uses \texttt{total\_CFs=1}.
\end{enumerate}
No explicit timeout/memory cap is set in EXP2 YAML; batch execution enforces one
evaluation worker per experiment process.

\subsubsection{Evaluation Metrics and Measurement Procedure}
Primary metrics are computed per instance in
\path{src/experiment/metrics_engine.py}:
\begin{enumerate}[leftmargin=*]
    \item \texttt{cost} (milliseconds),
    \item \texttt{sparsity} (active-feature ratio, threshold $\tau=10^{-4}$),
    \item \texttt{fidelity} (faithfulness correlation score),
    \item \texttt{faithfulness\_gap} (top-$k$ masking gap, $k=5$),
    \item \texttt{stability} (mean pairwise cosine similarity under perturbations).
\end{enumerate}
Faithfulness implementation (\path{src/metrics/faithfulness.py}):
\begin{align}
\Delta_k &= \left|f(x)-f\left(x_{\text{mask-top-}k}\right)\right|,\\
\rho &= \mathrm{corr}\left(|w_i|,\left|f(x)-f(x_{-i})\right|\right).
\end{align}
Mapping in pipeline: \texttt{fidelity}$\leftarrow \rho$ and
\texttt{faithfulness\_gap}$\leftarrow \Delta_k$.

Stability implementation (\path{src/metrics/stability.py}):
\begin{align}
x^{(t)} &= x + \epsilon^{(t)},\qquad \epsilon^{(t)}\sim\mathcal{N}(0,\sigma^2),\\
S &= \frac{2}{T(T-1)}\sum_{a<b}\cos\left(e^{(a)},e^{(b)}\right).
\end{align}
Secondary metrics: \texttt{cf\_sensitivity} only when
\texttt{counterfactual=true}; domain alignment implemented but disabled in EXP2.
Aggregation hierarchy is instance $\rightarrow$ run $\rightarrow$ block.

EEU is specified in EXP1 metric-design configs but not computable in current
EXP2 runtime due missing memory/model-call telemetry.
\texttt{[TO FILL: implement EEU instrumentation or keep EEU out of Paper A claims]}.

\subsubsection{Statistical Analysis Plan}
Artifact-qualified run inventory is:
\begin{enumerate}[leftmargin=*]
    \item planned runs: 300,
    \item present \texttt{results.json}: 250,
    \item analyzable runs: 233,
    \item excluded artifacts: 16 empty and 1 malformed JSON.
\end{enumerate}
Coverage imbalance is shown in Figure~\ref{fig:coverage}.

\begin{figure}[t]
\centering
\includegraphics[width=0.95\linewidth]{../../outputs/paper_a_figures/figures/fig_a2_coverage_heatmap.pdf}
\caption{Coverage map of the EXP2 robustness grid by model and explainer. Cell text reports files present out of 15 expected (\(5\) seeds \(\times\) \(3\) sample sizes), with invalid artifact counts where applicable.}
\label{fig:coverage}
\end{figure}

Global inference is restricted to 15 complete $(\text{model},N)$ blocks. SHAP--LIME
matched sets are 45 cells (\texttt{logreg/rf/xgb}) and 65 cells (all models).
Inferential stack (\path{scripts/run_exp2_statistical_analysis.py}):
\begin{enumerate}[leftmargin=*]
    \item Friedman omnibus test per metric.
    \item Nemenyi post-hoc localization per metric.
    \item Two-sided Wilcoxon signed-rank SHAP--LIME tests on matched cells.
\end{enumerate}
Multiplicity control: Nemenyi within metric; Holm-Bonferroni across the five
primary metrics for each inferential family (Friedman, Wilcoxon-45, Wilcoxon-65).
Uncertainty/effect reporting: mean, std, CV, $t$ and BCa bootstrap CIs, Kendall's
$W$, Cohen's $d_z$, and median paired differences.

\subsubsection{Robustness and Sensitivity Analyses}
Implemented robustness axes are seed heterogeneity (5 seeds), sampling-intensity
heterogeneity ($N=50/100/200$), and Gaussian perturbation stability.
Sensitivity controls:
\begin{enumerate}[leftmargin=*]
    \item Re-estimate SHAP--LIME contrasts on both 45-cell and 65-cell matched sets.
    \item Restrict global inference to block-complete contexts.
    \item Exclude malformed/empty artifacts without reconstruction.
\end{enumerate}
Robustness claim rule:
\begin{enumerate}[leftmargin=*]
    \item SHAP--LIME median-difference sign must match between 45-cell and 65-cell analyses.
    \item Holm-adjusted significance decision must not flip between the two analyses.
\end{enumerate}
Adversarial stress testing is not part of Paper A.
\texttt{[TO FILL: adversarial stress protocol with pre-registered deviation thresholds]}.

\subsubsection{Reproducibility Package}
Released package includes code, experiment YAMLs, frozen model/preprocessor
artifacts, run outputs, aggregate outputs, inferential exports, and figure
generation scripts.
Reference provenance anchors:
\begin{enumerate}[leftmargin=*]
    \item Host used for manuscript regeneration: Apple M3 Pro, macOS 26.3 (arm64), 11 logical CPUs, 18 GB RAM, Python 3.13.2.
    \item \texttt{sha256(requirements-frozen.txt)} = \texttt{b52d5a1f2e2edd5ada372ca66d18ef1447712fdd2c21f004d7e1f46a5ef9c6dc}.
    \item \path{outputs/batch_manifest.json} git hash \texttt{9fc70eb1e218a11f2fdd4bdb3aab3ea10a799a6f}.
\end{enumerate}
Minimal rerun sequence:
\begin{enumerate}[leftmargin=*]
    \item \texttt{python3 scripts/run\_batch\_experiments.py --config-dir configs/experiments/exp2\_scaled --output outputs/batch\_results.csv --parallel}
    \item \texttt{python3 scripts/check\_missing\_results.py}
    \item \texttt{python3 scripts/quick\_summary.py}
    \item \texttt{python3 scripts/full\_summary.py}
    \item \texttt{python3 scripts/generate\_paper\_a\_figures.py}
    \item \texttt{python3 scripts/run\_exp2\_statistical\_analysis.py}
\end{enumerate}
Archival status:
\texttt{[TO FILL: immutable DOI bundle for camera-ready reproducibility package]}.

\subsection{Inferential Implementation Details}
This subsection defines the inferential workflow for EXP2 and its reproducibility
support cohorts.

\subsubsection{Inferential Objective and RQ Linkage}
The statistical design implements Section~3.2:
\begin{enumerate}[leftmargin=*]
    \item RQ1 (method separation): test whether methods differ for each metric.
    \item RQ2 (trade-off stability): test whether method contrasts remain stable
    across model families and sampling intensities.
\end{enumerate}
All tests are executed on run/block-level units (Section~3.4), not on pooled
instance rows.

\subsubsection{Analysis Populations and Eligibility}
Inferential inputs are artifact-qualified EXP2 outputs:
\begin{enumerate}[leftmargin=*]
    \item Planned robustness grid: 300 runs.
    \item Present artifacts: 250 files.
    \item Analyzable artifacts: 233 (non-empty \texttt{instance\_evaluations}).
    \item Invalid artifacts: 1 malformed JSON and 16 empty files.
\end{enumerate}
Eligibility rules follow Section~3.4.5: parseable JSON, non-empty run metrics,
no imputation, and block-complete filtering for omnibus tests.

Resulting inferential sets are:
\begin{enumerate}[leftmargin=*]
    \item Global omnibus set: 15 complete $(\text{model},N)$ blocks $\times$ 4 methods.
    \item Primary SHAP--LIME paired set: 45 matched cells (\texttt{logreg/rf/xgb}).
    \item Sensitivity SHAP--LIME paired set: 65 matched cells (all model families).
    \item EXP1 reproducibility cohort: 9-seed repeated runs per core configuration.
\end{enumerate}

\subsubsection{Estimands, Aggregation, and Metric Directionality}
For metric $m$, run-level estimands are $y_{g,k,s,n,m}$ and block summaries are:
\[
\bar{y}_{k,b,m}=\frac{1}{|S_{k,b}|}\sum_{s\in S_{k,b}}y_{g,k,s,n,m},\quad b=(g,n).
\]
Primary metrics are fidelity, stability, sparsity, faithfulness gap, and cost.
Directionality is fixed pre-analysis:
\begin{enumerate}[leftmargin=*]
    \item Higher is better: fidelity, stability, faithfulness gap.
    \item Lower is better: sparsity, cost.
\end{enumerate}
Directionality is used for interpretation and normalized visual summaries (see
Figure X), while inferential tests operate on observed values per block.

\subsubsection{Global Omnibus Test (Friedman)}
For each metric:
\begin{enumerate}[leftmargin=*]
    \item Build block-by-method matrix (rows: complete $(\text{model},N)$ blocks;
    columns: methods).
    \item Apply Friedman test with methods as treatments and blocks as matched contexts.
    \item Evaluate significance at $\alpha=0.05$ (Table X).
\end{enumerate}
Friedman is used because the design is repeated-measures and metric
distributions are not assumed Gaussian.

\subsubsection{Post-hoc Localization (Nemenyi)}
If Friedman rejects the null for metric $m$:
\begin{enumerate}[leftmargin=*]
    \item Run Nemenyi pairwise post-hoc comparisons on the same block matrix.
    \item Report pairwise adjusted p-values for all method pairs.
    \item Use $\alpha=0.05$ for within-metric family-wise control.
\end{enumerate}
Across-metric multiplicity is controlled via Holm-Bonferroni adjustment over the
five primary metrics, applied separately to each inferential family.

\subsubsection{Focused SHAP--LIME Paired Test (Wilcoxon)}
For each primary metric:
\begin{enumerate}[leftmargin=*]
    \item Construct matched SHAP/LIME pairs on identical $(\text{model},\text{seed},N)$ cells.
    \item Apply two-sided Wilcoxon signed-rank test.
    \item Report paired descriptive summaries and p-values.
\end{enumerate}
Primary analysis uses the 45-cell restricted set; 65-cell all-model analysis is
used as sensitivity support.

\subsubsection{Uncertainty Quantification}
Reported uncertainty statistics include mean, standard deviation, and CV. Where
configured, confidence intervals are computed via
\path{src/analysis/confidence.py}:
\begin{enumerate}[leftmargin=*]
    \item $t$-distribution CIs (\texttt{compute\_t\_ci}).
    \item Bootstrap CIs (\texttt{compute\_bootstrap\_ci}), BCa method with
    10,000 resamples and seed 42; percentile fallback on failure.
\end{enumerate}

\subsubsection{Effect Size and Practical Significance}
Effect-size outputs are generated alongside p-values:
\begin{enumerate}[leftmargin=*]
    \item Kendall's $W$ for Friedman tests.
    \item Cohen's $d_z$ for paired Wilcoxon SHAP--LIME contrasts.
    \item Median paired differences for location-scale interpretation.
\end{enumerate}
These are exported by the EXP2 statistical driver for direct appendix use.

\subsubsection{Sensitivity and Missingness Controls}
Sensitivity controls in the analysis layer:
\begin{enumerate}[leftmargin=*]
    \item Re-estimate SHAP--LIME contrasts on restricted (45) and all-model (65) matched sets.
    \item Restrict global inference to block-complete contexts.
    \item Exclude malformed/empty artifacts without synthetic reconstruction.
\end{enumerate}
Dispersion is quantified via std/CV/CI. No threshold-based practical-equivalence
rule is enforced in the current protocol; robustness interpretation therefore
relies on direction consistency, effect-size magnitude, and uncertainty bounds.

\subsubsection{Reproducible Analysis Execution}
Current reproducible analysis stack:
\begin{enumerate}[leftmargin=*]
    \item Batch execution and artifact generation.
    \item Missing/malformed audit (\texttt{scripts/check\_missing\_results.py}).
    \item Summary tables (\texttt{scripts/quick\_summary.py}, \texttt{scripts/full\_summary.py}).
    \item Figure regeneration (\texttt{scripts/generate\_paper\_a\_figures.py}).
    \item Deterministic EXP2 inference export:
    \texttt{python3 scripts/run\_exp2\_statistical\_analysis.py}.
\end{enumerate}
The EXP2 driver materializes Friedman, Nemenyi, Wilcoxon, multiplicity-adjusted
p-values, effect-size tables, and matched-cell inventories in
\path{outputs/analysis/paper_a_exp2_stats/}.

\subsection{Reproducibility Protocol}
\subsubsection{Scope and Methodological Role}
This protocol defines how Section~4 quantitative claims are regenerated from
repository artifacts. It links Sections~3.6--3.7 to a deterministic execution
contract for RQ1/RQ2 (see Table~X).

\subsubsection{Definitions, Notation, and Assumptions}
Run key is $r=(g,k,s,n)$ with model family $g$, explainer $k$, seed $s$, and
sampling intensity $n$.
Reproducibility state is decomposed into:
\begin{enumerate}[leftmargin=*]
    \item $\mathcal{C}$: configuration state (YAMLs, seed grid, CLI invocations),
    \item $\mathcal{A}$: artifact state (\texttt{results.json}, \texttt{metrics.csv}, aggregate exports),
    \item $\mathcal{I}$: inferential state (\path{outputs/analysis/paper_a_exp2_stats/}),
    \item $\mathcal{V}$: variance state (\texttt{reproducibility\_report.csv}).
\end{enumerate}
Assumptions: EXP1 model artifacts are frozen during EXP2; config files are
immutable during one run; only artifact-qualified runs enter inference.

\subsubsection{Executable Reproducibility Procedure}
\begin{enumerate}[leftmargin=*]
    \item Freeze protocol anchors:
    \begin{enumerate}[leftmargin=*]
        \item config roots (\path{configs/experiments/exp2_scaled}, \path{configs/experiments/exp2_comparative}),
        \item dependency hash: \texttt{sha256(requirements-frozen.txt)}=
        \texttt{b52d5a1f2e2edd5ada372ca66d18ef1447712fdd2c21f004d7e1f46a5ef9c6dc},
        \item provenance hash in \path{outputs/batch_manifest.json}.
    \end{enumerate}
    \item Execute EXP2 batch generation.
    \item Audit missing/malformed/empty artifacts.
    \item Regenerate deterministic inference exports via
    \texttt{scripts/run\_exp2\_statistical\_analysis.py}.
    \item Regenerate summaries/figures and repeated-run reproducibility outputs.
\end{enumerate}

\subsubsection{Quality Checks and Acceptance Criteria}
A reproduction pass is claim-eligible only if:
\begin{enumerate}[leftmargin=*]
    \item seed/model/explainer grids match declared manifests;
    \item malformed and empty artifacts are explicitly enumerated;
    \item omnibus inference uses only complete $(\text{model},N)$ blocks;
    \item Section~4 claims map to versioned files under analysis/figure outputs;
    \item repeated-run variance diagnostics are present.
\end{enumerate}

\subsubsection{Environment and Compute Disclosure}
Reference regeneration host: Apple M3 Pro, macOS 26.3 (arm64), 11 logical CPUs,
18 GB RAM, Python 3.13.2. This is a provenance anchor, not an execution
requirement. \texttt{[TO FILL: multi-hardware normalization protocol if runtime claims are generalized across platforms]}.

\subsection{Validity Threats and Mitigations}
\subsubsection{Scope}
Threat analysis is organized by validity class to separate implemented controls
from residual uncertainties.

\subsubsection{Internal Validity}
\begin{enumerate}[leftmargin=*]
    \item Threat: incomplete EXP2 execution coverage (300 planned, 250 present, 233 analyzable).
    \item Mitigation: explicit artifact audit, block-complete filtering, matched-cell filtering.
    \item Residual risk: potential non-random missingness by method--model cell.
\end{enumerate}

\subsubsection{Construct Validity}
\begin{enumerate}[leftmargin=*]
    \item Threat: ambiguity in metric naming/interpretation.
    \item Mitigation: code-level operationalization and explicit mapping
    (\texttt{fidelity}$\leftarrow$faithfulness correlation,
    \texttt{faithfulness\_gap}$\leftarrow$top-$k$ masking shift).
    \item Residual risk: EEU not computed in current EXP2 runtime; cost is wall-clock only.
\end{enumerate}

\subsubsection{Statistical Conclusion Validity}
\begin{enumerate}[leftmargin=*]
    \item Threat: inflated Type-I error and pseudo-replication.
    \item Mitigation: run/block analysis units, Friedman+Nemenyi+Wilcoxon stack, Holm adjustments.
    \item Residual risk: no pre-registered practical-equivalence threshold.
\end{enumerate}

\subsubsection{External Validity}
\begin{enumerate}[leftmargin=*]
    \item Threat: single benchmark dataset/task family.
    \item Mitigation: diversity across five model families, four explainers, five seeds, and three sample sizes.
    \item Residual risk: generalization to text/vision/time-series is not established in Paper A.
\end{enumerate}

\subsubsection{Implementation Validity}
\begin{enumerate}[leftmargin=*]
    \item Threat: YAML-exposed parameters not fully bound at runtime (Anchors threshold, DiCE CF count).
    \item Mitigation: explicit disclosure of effective runtime settings.
    \item Residual risk: harmonization could shift absolute method levels.
\end{enumerate}
\texttt{[TO FILL: camera-ready harmonization diff report between YAML-exposed and runtime-bound parameters]}.

\subsection{Framework Operation Method (FOM-7): Executable Protocol}
\subsubsection{Formal Definition and Role}
FOM-7 is the operational method contribution of Paper A. It is defined as:
\[
\mathcal{M}_{\mathrm{FOM7}}=\{(S_i,A_i,G_i)\}_{i=1}^{7},
\]
with stage $S_i$, required artifacts $A_i$, and gate condition $G_i$.
Progression is allowed only if $G_i=\mathrm{pass}$.

\subsubsection{Stage Protocol and Gate Logic}
Algorithm 1 (FOM-7) maps protocol execution to stage-gated evidence production:
\begin{enumerate}[leftmargin=*]
    \item \textbf{S1 Protocol specification}: freeze study design and enumerate the run grid.
    \item \textbf{S2 Controlled execution}: run batch jobs and persist per-run artifacts plus execution manifest.
    \item \textbf{S3 Integrity audit}: enumerate missing, empty, and malformed artifacts.
    \item \textbf{S4 Harmonization}: aggregate retained outputs into analysis-ready metric tables.
    \item \textbf{S5 Statistical inference}: run Friedman, Nemenyi, and paired Wilcoxon tests.
    \item \textbf{S6 Reproducibility characterization}: estimate cross-seed dispersion and CV.
    \item \textbf{S7 Claim-ready reporting}: map manuscript claims to versioned artifacts and caveat ledger.
\end{enumerate}
Progression from stage $S_i$ to stage $S_{i+1}$ is allowed only when the corresponding gate $G_i$ is satisfied.

\begin{figure}[t]
\centering
\includegraphics[width=0.98\linewidth]{../../outputs/paper_a_figures/figures/fig_a1_fom7_flow.pdf}
\caption{Operational view of FOM-7. The workflow is stage-gated: each stage must produce auditable artifacts before progression to the next stage.}
\label{fig:fom7-flow}
\end{figure}
Figure~\ref{fig:fom7-flow} provides the operational view of FOM-7.

\subsubsection{Reference Implementation Procedure}
\begin{enumerate}[leftmargin=*]
    \item \texttt{python3 scripts/run\_batch\_experiments.py}\\
    \texttt{--config-dir configs/experiments/exp2\_scaled}\\
    \texttt{--output outputs/batch\_results.csv --parallel}
    \item \texttt{python3 scripts/check\_missing\_results.py}
    \item \texttt{python3 scripts/quick\_summary.py}
    \item \texttt{python3 scripts/full\_summary.py}
    \item \texttt{python3 scripts/run\_reproducibility\_study.py}\\
    \texttt{--config-dir configs/experiments --pattern "exp1\_adult\_*.yaml"}\\
    \texttt{--output-dir experiments/exp1\_adult/reproducibility}
    \item \texttt{python3 scripts/generate\_paper\_a\_figures.py}
    \item \texttt{python3 scripts/run\_exp2\_statistical\_analysis.py}
\end{enumerate}
If any gate fails, downstream inferential outputs are not claim-eligible and the
affected comparison is downgraded to descriptive-only reporting.

\section{Results}

\subsection{Global Trade-offs (EXP2 robustness)}
\begin{figure}[t]
\centering
\includegraphics[width=0.90\linewidth]{../../outputs/paper_a_figures/figures/fig_a3_method_radar.pdf}
\caption{Normalized method profile across core metrics (higher is better after direction-aware normalization). SHAP dominates fidelity/stability, while LIME and DiCE dominate efficiency/sparsity dimensions, respectively.}
\label{fig:method-radar}
\end{figure}
Figure~\ref{fig:method-radar} provides a direction-normalized profile view. Table~\ref{tab:global-tradeoffs} reports absolute means over 15 complete model-size blocks.

\begin{table}[h]
\centering
\caption{Method-level means over complete model-size blocks.}
\begin{tabular}{lrrrrr}
\toprule
Method & Fidelity & Stability & Sparsity & Faithfulness Gap & Time (ms) \\
\midrule
SHAP    & 0.8176 & 0.7377 & 0.2264 & 0.4474 & 685220.23 \\
LIME    & 0.5602 & 0.0144 & 0.0846 & 0.3342 & 3660.68 \\
Anchors & 0.3853 & 0.0006 & 0.0928 & 0.2382 & 25326.92 \\
DiCE    & 0.1716 & 0.3602 & 0.0164 & 0.0988 & 16306.50 \\
\bottomrule
\end{tabular}
\label{tab:global-tradeoffs}
\end{table}

\begin{figure}[t]
\centering
\includegraphics[width=0.88\linewidth]{../../outputs/paper_a_figures/figures/fig_a4_fidelity_cost_tradeoff.pdf}
\caption{Fidelity--cost trade-off frontier using method means over complete blocks (error bars: block-level standard deviation). The x-axis is logarithmic due strong runtime heterogeneity across methods.}
\label{fig:fidelity-cost}
\end{figure}
Figure~\ref{fig:fidelity-cost} isolates the fidelity--cost frontier and its dispersion across complete blocks.

\subsection{Statistical Significance}
\begin{table}[h]
\centering
\caption{Friedman tests across methods.}
\begin{tabular}{lrrl}
\toprule
Metric & Statistic & p-value & Significant \\
\midrule
Fidelity         & 42.12 & 3.78e-09 & Yes \\
Stability        & 43.88 & 1.60e-09 & Yes \\
Sparsity         & 35.64 & 8.92e-08 & Yes \\
Faithfulness Gap & 45.00 & 9.25e-10 & Yes \\
Cost             & 27.72 & 4.16e-06 & Yes \\
\bottomrule
\end{tabular}
\label{tab:friedman}
\end{table}

\subsection{Focused SHAP vs LIME Comparison}
\begin{table}[h]
\centering
\caption{Paired Wilcoxon comparison on 45 matched configurations (\texttt{logreg/rf/xgb}).}
\begin{tabular}{lrrr}
\toprule
Metric & SHAP Mean & LIME Mean & Wilcoxon p-value \\
\midrule
Fidelity         & 0.8112 & 0.5556 & 5.68e-14 \\
Stability        & 0.8013 & 0.0154 & 5.68e-14 \\
Sparsity         & 0.3156 & 0.0845 & 5.68e-14 \\
Faithfulness Gap & 0.3924 & 0.3408 & 5.68e-14 \\
Cost (ms)        & 1258.72 & 210.45 & 2.29e-06 \\
\bottomrule
\end{tabular}
\label{tab:shap-lime}
\end{table}

\subsection{Reproducibility Snapshot (EXP1 repeated runs)}
Fidelity CV $\leq 0.0634$, stability CV $\leq 0.0826$, sparsity CV $\leq 0.0441$, and faithfulness-gap CV $\leq 0.0358$ across baseline configurations. Cost CV reaches 0.225 for SHAP variants.

\begin{figure}[t]
\centering
\includegraphics[width=0.92\linewidth]{../../outputs/paper_a_figures/figures/fig_a5_reproducibility_cv.pdf}
\caption{Reproducibility heatmap for EXP1 repeated runs (9 seeds per core configuration). Lower CV indicates stronger run-to-run stability.}
\label{fig:repro-cv}
\end{figure}
Figure~\ref{fig:repro-cv} highlights that cost variance dominates quality-metric variance in repeated runs.

\section{Validity and Reporting Caveats}
\begin{enumerate}[leftmargin=*]
    \item EXP2 robustness is incomplete (50 missing runs) with one malformed result file.
    \item Runtime comparisons are sensitive to implementation details and hardware.
    \item Semantic evaluation is out of scope for Paper A.
\end{enumerate}

\section{JMLR Track Positioning}
\begin{enumerate}[leftmargin=*]
    \item Reproducible benchmark framework with artifact lineage and failure accounting.
    \item Multi-metric protocol with non-parametric inferential testing.
    \item Reusable operation method (FOM-7) for benchmark execution and audit standardization.
    \item Clearly scoped quantitative benchmark; semantic track deferred.
\end{enumerate}

\section{Conclusion}
The evidence supports a clear trade-off frontier: SHAP leads fidelity/stability, LIME leads runtime practicality, and DiCE leads sparsity. In addition to empirical findings, Paper A contributes FOM-7 as a concrete operating protocol for running, auditing, and reporting XAI benchmark studies.

\acks{Funding and competing-interest disclosures are to be finalized in the submission package. This draft was prepared from repository artifacts dated February 2026.}

\end{document}
